\begin{helper}[GnpVertexIncidentEdgeGeneratingFunction]
Fix \(m\), \(p,z\in\mathbb R\), and \(r\in\operatorname{Fin}(m)\).  Let
\[
E_r=\{(a,b):a,b\in\operatorname{Fin}(m),\ a<b,\ a=r\text{ or }b=r\}
\]
be the represented edge coordinates incident to \(r\).  Under the
\(\noderef{GnpGraphWeight}\) product graph weight,
\[
\sum_G
\left(\prod_{e=(a,b)\in E_r}
  \begin{cases}
  z,&G.Adj(a,b),\\
  1,&\neg G.Adj(a,b)
  \end{cases}\right)
\noderef{GnpGraphWeight}(p,G)
=(1-p+pz)^{|E_r|}.
\]
\end{helper}
\begin{proof}
Expand \(\noderef{GnpGraphWeight}\) as the product over all represented
unordered edge coordinates.  Summing over graphs is the same as summing over
Boolean assignments to those coordinates.  For a coordinate not incident to
\(r\), the extra product contributes \(1\), and the graph-weight coordinate
sum is \(p+(1-p)=1\).  For a coordinate incident to \(r\), the absent choice
contributes \(1\cdot(1-p)\), while the present choice contributes \(z\cdot p\),
so the coordinate sum is \(1-p+pz\).  The finite product-sum factorization is
the same coordinatewise collapse encoded by \(\noderef{ProductWeightSumOne}\).
Multiplying the incident-coordinate factors gives \((1-p+pz)^{|E_r|}\), and
all nonincident factors are \(1\).
\end{proof}
